% Exposición oficial — Plan de Tesis Grupo 9.2 (Beamer PDF, 22 slides)
% Compilar: make pdf  →  ../Exposicion_Plan_Tesis_Grupo9.2-beamer.pdf
\documentclass[aspectratio=169,11pt,spanish]{beamer}

% Metadatos del Plan de Tesis — Grupo 9.2 (Anexo III / Guía N°01)
\newcommand{\tituloPlan}{%
  Alineación entre contratos OpenAPI bancarios y modelos de referencia BIAN
  mediante análisis estructural y similitud semántica}
\newcommand{\tituloPlanCaps}{%
  ALINEACIÓN ENTRE CONTRATOS OPENAPI BANCARIOS Y MODELOS DE REFERENCIA BIAN\\
  MEDIANTE ANÁLISIS ESTRUCTURAL Y SIMILITUD SEMÁNTICA}

\newcommand{\grupo}{Grupo 9.2}
\newcommand{\anioPlan}{2026}
\newcommand{\ciudad}{Lima, Perú}
\newcommand{\ciudadCaps}{LIMA, PERÚ}

\newcommand{\facultadCaps}{%
  UNIDAD DE POSGRADO FACULTAD DE INGENIERÍA DE SISTEMAS Y COMPUTACIÓN\\
  / ESCUELA DE POSGRADO}
\newcommand{\mencionGrado}{INTELIGENCIA ARTIFICIAL}

\newcommand{\autorAlex}{Alex Mancilla Antay}
\newcommand{\autorJack}{Jack Paitan Cano}
\newcommand{\autorAlexCaps}{ALEX MANCILLA ANTAY}
\newcommand{\autorJackCaps}{JACK PAITAN CANO}

\newcommand{\asesorPlan}{Por designar}
\newcommand{\asesorPlanCaps}{POR DESIGNAR}

\newcommand{\areaTematica}{Inteligencia Artificial --- alineación de APIs bancarias OpenAPI con modelos BIAN}
\newcommand{\institucionDesarrollo}{Universidad Nacional de Ingeniería}
\newcommand{\duracionMeses}{6}
\newcommand{\duracionFechas}{Marzo 2026 --- Junio 2026}

\newcommand{\palabrasClaveES}{%
  OpenAPI; BIAN; AlignmentScore; schema matching; API bancaria; alineación estructural}
\newcommand{\palabrasClaveEN}{%
  OpenAPI; BIAN; AlignmentScore; schema matching; banking API; structural alignment}

% Informe Guía 1 — URL Google Drive (última hoja). Sustituir antes de entregar.
\newcommand{\urlDriveEntrega}{https://drive.google.com/drive/folders/1Kuc3eIuBjt_EXaPislf0hqwSEX-kIdSI?usp=drive_link}

% Estilo exposición Beamer — Grupo 9.2 (16:9, español, Helvetica)
\usepackage[utf8]{inputenc}
\usepackage[T1]{fontenc}
\usepackage[spanish,es-tabla]{babel}
\usepackage{graphicx}
\usepackage{adjustbox}
\usepackage{array}
\usepackage{colortbl}
\usepackage{booktabs}
\usepackage{ragged2e}
\usepackage{tabularx}
\usepackage{tcolorbox}
\tcbuselibrary{skins}
\usepackage{tikz}
\usepackage{xurl}

\usetheme{Madrid}
\usecolortheme{whale}
\usefonttheme{professionalfonts}
\usepackage[scaled=0.92]{helvet}
\renewcommand{\familydefault}{\sfdefault}

\definecolor{uninavy}{RGB}{27,54,93}
\definecolor{uniaccent}{RGB}{232,163,23}
\definecolor{unimuted}{RGB}{99,110,114}
\definecolor{unidanger}{RGB}{192,57,43}
\definecolor{tablehead}{RGB}{27,54,93}
\definecolor{tablealt}{RGB}{238,242,247}
\definecolor{tablerule}{RGB}{206,214,224}
\definecolor{tableframe}{RGB}{27,54,93}

\setbeamercolor{structure}{fg=uninavy}
\setbeamercolor{frametitle}{bg=uninavy,fg=white}
\setbeamercolor{title}{fg=white,bg=uninavy}
\setbeamercolor{block title}{bg=uninavy,fg=white}
\setbeamercolor{block body}{bg=uninavy!5,fg=black}
\setbeamercolor{block title alerted}{bg=unidanger,fg=white}
\setbeamercolor{block body alerted}{bg=unidanger!8,fg=black}
\setbeamercolor{block title example}{bg=uniaccent,fg=uninavy}
\setbeamercolor{block body example}{bg=uniaccent!14,fg=black}
\setbeamercolor{footline}{fg=unimuted,bg=white}

\setbeamertemplate{navigation symbols}{}
\setbeamertemplate{blocks}[rounded][shadow=false]
\setbeamersize{text margin left=11mm, text margin right=11mm}

\setbeamertemplate{footline}{%
  \leavevmode%
  \hbox{%
    \begin{beamercolorbox}[wd=.72\paperwidth,ht=2.5ex,dp=1.2ex,leftskip=1em]{footline}%
      \grupo{} · Plan de Tesis · UNI · Jul 2026%
    \end{beamercolorbox}%
    \begin{beamercolorbox}[wd=.28\paperwidth,ht=2.5ex,dp=1.2ex,rightskip=1em]{footline}%
      \insertframenumber{} / \inserttotalframenumber%
    \end{beamercolorbox}%
  }%
}

\setbeamertemplate{itemize item}{\color{uniaccent}\small$\blacktriangleright$}
\setbeamertemplate{itemize subitem}{\color{uninavy!70}\scriptsize$\circ$}
\setbeamerfont{normal text}{size=\large}
\setbeamerfont{item}{size=\large}
\setbeamerfont{frametitle}{size=\large,series=\bfseries}
\setbeamerfont{block title}{size=\normalsize,series=\bfseries}
\setbeamerfont{block body}{size=\large}

\graphicspath{%
  {ppt-images/}%
  {../../figuras/cap01/}%
  {../../figuras/cap02/}%
  {../../figuras/cap03/}%
  {../../figuras/modelo-alineacion/}%
}

\newcommand{\figplan}[2][]{%
  \vspace{-0.15em}%
  \centering
  \adjustbox{max width=0.94\linewidth, max height=0.68\textheight, center}{%
    \includegraphics[#1]{#2}%
  }%
}

\newcommand{\figplanlarge}[2][]{%
  \vspace{-0.35em}%
  \centering
  \adjustbox{max width=0.98\linewidth, max height=0.80\textheight, center}{%
    \includegraphics[#1]{#2}%
  }%
}

% Figura ontología (slide 4 · no oral) — máximo espacio vertical
\newcommand{\figplanontology}[2][]{%
  \vspace{-0.45em}%
  \centering
  \adjustbox{max width=\linewidth, max height=0.83\textheight, center}{%
    \includegraphics[#1]{#2}%
  }%
}

\newcommand{\figplancompact}[2][]{%
  \vspace{-0.15em}%
  \centering
  \adjustbox{max width=0.92\linewidth, max height=0.50\textheight, center}{%
    \includegraphics[#1]{#2}%
  }%
}

\newcommand{\figppttwo}[2]{%
  \vspace{-0.2em}%
  \centering
  \adjustbox{max width=0.96\linewidth, max height=0.34\textheight, center}{%
    \includegraphics{#1}%
  }\\[0.12em]
  \adjustbox{max width=0.96\linewidth, max height=0.34\textheight, center}{%
    \includegraphics{#2}%
  }%
}

\newcommand{\figcaption}[1]{%
  \vspace{0.2em}%
  \parbox{0.94\linewidth}{\centering\large\bfseries\textcolor{uninavy}{#1}}%
}

% --- Tablas ---
\newcolumntype{L}[1]{>{\RaggedRight\arraybackslash}p{#1}}
\newcolumntype{Y}{>{\RaggedRight\arraybackslash}X}
\newcommand{\tabhead}[1]{\color{white}\textbf{#1}}
\newcommand{\tabphase}[1]{\textcolor{uninavy}{\textbf{#1}}}

\newcommand{\plantablesetup}{%
  \setlength{\tabcolsep}{5.5pt}%
  \renewcommand{\arraystretch}{1.14}%
  \arrayrulecolor{tablerule}%
}

% Marco ancho completo · esquinas redondeadas · franja superior dorada
\newcommand{\beamertable}[1]{%
  \par\vspace{0.1em}%
  \begin{tcolorbox}[
    enhanced,
    arc=4pt,
    boxrule=0.55pt,
    colframe=tableframe,
    colback=white,
    width=\linewidth,
    boxsep=0pt,
    left=7pt, right=7pt, top=5pt, bottom=5pt,
    borderline north={1.05pt}{0pt}{uniaccent},
  ]
  {\fontsize{13.5}{16.2}\selectfont\plantablesetup #1}
  \end{tcolorbox}
  \par\vspace{0.2em}%
}

\newcommand{\beamertablecompact}[1]{%
  \par\vspace{0.03em}%
  \begin{tcolorbox}[
    enhanced,
    arc=3pt,
    boxrule=0.55pt,
    colframe=tableframe,
    colback=white,
    width=\linewidth,
    boxsep=0pt,
    left=4pt, right=4pt, top=3pt, bottom=3pt,
    borderline north={1.05pt}{0pt}{uniaccent},
  ]
  {%
    \fontsize{7.75}{9.4}\selectfont
    \plantablesetup
    \setlength{\tabcolsep}{3pt}%
    \renewcommand{\arraystretch}{1.0}%
    #1%
  }
  \end{tcolorbox}
  \par\vspace{0.1em}%
}

\newcommand{\tablesectiontitle}[1]{%
  \par\vspace{0.22em}%
  \noindent
  {\color{uniaccent}\rule{0.045\linewidth}{2pt}\hspace{0.35em}%
   \small\bfseries\textcolor{uninavy}{#1}}%
  \par\vspace{0.08em}%
}

% Bloque recordatorio (acento dorado)
\newenvironment{planreminder}[1]{%
  \begin{exampleblock}{#1}%
}{%
  \end{exampleblock}%
}

% Lista en panel suave (sin título)
\newenvironment{planpanel}{%
  \begin{beamercolorbox}[wd=0.98\linewidth,sep=0.55em,rounded=true,shadow=false]{block body}%
}{%
  \end{beamercolorbox}%
}

\newcommand{\keyterm}[1]{\textbf{\textcolor{uninavy}{#1}}}

% Punto rojo esquina superior derecha — slide no oral
\newcommand{\marcadorNoExponer}{%
  \begin{tikzpicture}[remember picture,overlay]
    \fill[unidanger] ([xshift=-0.65cm,yshift=-0.55cm]current page.north east) circle (6pt);
  \end{tikzpicture}%
}


\title[Plan de Tesis Grupo 9.2]{\tituloPlan}
\author{\autorAlex \and \autorJack}
\institute{\institucionDesarrollo · \grupo}
\date{\ciudad · \anioPlan}

\begin{document}

% --- 1 Portada ---
\begin{frame}[plain]
  \centering
  \vspace{0.1cm}
  {\large\bfseries UNIVERSIDAD NACIONAL DE INGENIERÍA\par}
  \vspace{0.2cm}
  {\small\bfseries Plan de Tesis · \grupo\par}
  \vspace{0.45cm}
  {\small\bfseries ``\tituloPlan''\par}
  \vspace{0.55cm}
  {\small Para obtener el grado de Maestro en Ciencias\\[0.12em]
  con mención en \mencionGrado\par}
  \vspace{0.5cm}
  {\small\bfseries \autorAlex\par}
  {\small\bfseries \autorJack\par}
  \vspace{0.35cm}
  {\small Investigación tecnológica · Design Science\par}
  {\small IEEE: Knowledge Representation · ACM: Semantic networks\par}
  \vspace{0.3cm}
  {\small \ciudad · \anioPlan\par}
  \note{[0:00--0:45] Jack. Título + tipo investigación + clasificación ACM/IEEE.}
\end{frame}

% --- 2 Fuentes (NO EXPONER) ---
\begin{frame}{Fuentes bibliográficas y gestión documental}
  \marcadorNoExponer
  \vspace{0.35em}
  \begin{planpanel}
  \begin{itemize}
    \item Link de los documentos:\\[0.15em]
      {\small\url{https://drive.google.com/drive/folders/1bNOd0ZvD_2-8HIaKEGm6qX1XRMpnxxab}}
    \item Link al Zotero Web:\\[0.15em]
      {\small\url{https://www.zotero.org/groups/6622247/proyectotesisgrupo9.2/library}}
  \end{itemize}
  \end{planpanel}
  \note{NO EXPONER. Saltar en clase (1 $\rightarrow$ 5).}
\end{frame}

% --- 3 Intro ---
\begin{frame}{Introducción al tema (§1.1)}
  \marcadorNoExponer
  \begin{planpanel}
  \small
  La industria bancaria utiliza contratos OpenAPI para documentar e integrar APIs entre canales digitales, fintechs y socios de negocio. Paralelamente, BIAN publica contratos OpenAPI de referencia para estandarizar las capacidades bancarias y promover la interoperabilidad entre entidades (Bhatia, 2022; Farzi, 2021).

  \vspace{0.35em}
  Sin embargo, actualmente no existe un procedimiento objetivo que permita medir el grado de alineación entre un contrato OpenAPI implementado por un banco y el contrato OpenAPI de referencia propuesto por BIAN para un mismo Service Domain.

  \vspace{0.35em}
  En respuesta a esta necesidad, la presente investigación propone un modelo de medición de la alineación basado en tres dimensiones: SemanticScore, StructuralScore y CoverageScore, integradas en un AlignmentScore, con el propósito de cuantificar de manera objetiva y reproducible el nivel de alineación entre ambos contratos OpenAPI.
  \end{planpanel}
  \note{NO EXPONER. Saltar en clase (1 $\rightarrow$ 5).}
\end{frame}

% --- 4 Ontología (NO EXPONER) ---
\begin{frame}{Ontología del plan --- conceptos relacionados}
  \marcadorNoExponer
  \figplanontology{Figura_02_Objeto_Unidad_Analisis}
  \vspace{0.05em}
  {\normalsize\bfseries\textcolor{uninavy}{Ontología interna --- objeto · UA · producto · referencia BIAN}}
  \note{NO EXPONER. Figura de respaldo (informe Cap.\ 1). Saltar en clase (1 $\rightarrow$ 5).}
\end{frame}

% --- 4 Problemática situación ---
\begin{frame}{Problemática --- situación real: dónde y actores (§1.2)}
  \figplan{slide05_img01}
  \figcaption{Situación real --- sector bancario, actores, open banking (CMA, 2016)}
  \note{[2:20--3:10] Jack. Situación real $\neq$ objeto $\neq$ problema tecnológico.}
\end{frame}

% --- 5 Síntomas ---
\begin{frame}{Problemática --- síntomas, dolor y estadísticas (§1.2)}
  \figppttwo{slide06_img01}{slide06_img02}
  \figcaption{Síntomas: TSB (2018) · fragmentación · Casas et al.\ (2021) 68\,\%}
  \note{[3:10--4:05] Jack. Tener PDF Casas listo (PI-13). Señalar cifra 68\,\%.}
\end{frame}

% --- 6 Objeto subyacente ---
\begin{frame}{Objeto de estudio --- situación subyacente (§1.4)}
  \figplan{slide07_img01}
  \figcaption{Situación subyacente --- contrato OpenAPI $\leftrightarrow$ Service Domain BIAN}
  \note{[4:05--4:50] Jack. BIAN = referencia, no objeto.}
\end{frame}

% --- 7 Objeto real ---
\begin{frame}{Objeto de estudio real --- taxonomía y ciclo de vida}
  \figplancompact{slide08_img01}
  \figcaption{Objeto real: taxonomía OpenAPI 3.x · ciclo de vida del contrato}
  \vspace{0.2em}
  \begin{planpanel}
  \begin{itemize}
    \item \keyterm{Entidad real:} contrato OpenAPI publicado por un banco
    \item \keyterm{Taxonomía:} OpenAPI 3.x · pagos / cuentas / préstamos · versiones
    \item \keyterm{Ciclo de vida:} diseño $\rightarrow$ revisión $\rightarrow$ publicación $\rightarrow$ consumo
  \end{itemize}
  \end{planpanel}
  \note{[4:50--5:30] Jack. Ciclo de vida del contrato, no del proceso bancario completo.}
\end{frame}

% --- 8 Instrumento ---
\begin{frame}{Instrumento y objeto modelado}
  \figplan{slide09_img01}
  \figcaption{Instrumento: parser OpenAPI + extracto BIAN $\rightarrow$ representaciones normalizadas}
  \note{[5:30--6:15] Alex. Instrumento = procedimiento, no encuesta ni IoT.}
\end{frame}

% --- 9 Casos A/B ---
\begin{frame}{Instrumento y objeto modelado --- casos e instancias}
  \figplan{slide10_img01}
  \figcaption{Casos A (S) y B (C) · clasificación Alta / Media / Baja / Nula}
  \note{[6:15--7:00] Alex. Ejemplos Anexo F.}
\end{frame}

% --- 10 Problema tecnológico ---
\begin{frame}{Problema tecnológico y preguntas P1--P4 (§1.3)}
  \begin{planreminder}{Recordatorio}
    \keyterm{Problemática del sector} $\neq$ \keyterm{problema tecnológico} del artefacto a diseñar
  \end{planreminder}
  \vspace{0.15em}
  \begin{planpanel}
  \begin{itemize}
    \item \keyterm{P1} matching estructural · \keyterm{P2} similitud semántica
    \item \keyterm{P3} cobertura · \keyterm{P4} score integrado (AlignmentScore)
    \item Métodos \keyterm{insuficientes} / no adaptados a OpenAPI$\leftrightarrow$BIAN (Shvaiko, 2005; Casas, 2021)
    \item Artefacto: procedimiento + modelo S/E/C $\rightarrow$ clasificación en [0,1]
  \end{itemize}
  \end{planpanel}
  \note{[7:00--7:50] Alex. Decir «insuficiente» en voz alta (PI-12).}
\end{frame}

% --- 11 OG ---
\begin{frame}{Objetivo general (técnico) --- insumo y variable dependiente}
  \begin{planpanel}
  \begin{itemize}
    \item \keyterm{Insumo:} contrato OpenAPI + Service Domain BIAN emparejado
    \item \keyterm{OG:} diseñar modelo S/E/C y procedimiento reproducible de alineación
    \item \keyterm{VD del artefacto:} AlignmentScore $\in$ [0,1] + clasificación
    \item \keyterm{Objetivo superior:} cuantificar coherencia antes del despliegue (consecuencia, no OG)
  \end{itemize}
  \end{planpanel}
  \note{[7:50--8:25] Alex. Separar OG técnico vs objetivo superior.}
\end{frame}

% --- 13 OE ---
\begin{frame}{Objetivos específicos OE1--OE4 (entrada $\rightarrow$ producto)}
  \figplanlarge{Figura_04_Cadena_Objetivos}
  \figcaption{Cada OE es hito con producto medible --- no «programar» ni «experimentar»}
  \note{[8:25--9:15] Alex. OE1 normaliza · OE2$\rightarrow$E · OE3$\rightarrow$S · OE4$\rightarrow$C + AlignmentScore.}
\end{frame}

% --- 13 VI ---
\begin{frame}{Variables independientes --- factores del investigador (Tabla 5)}
  \beamertable{%
  \begin{tabular}{@{}L{0.07\linewidth}L{0.29\linewidth}L{0.57\linewidth}@{}}
    \rowcolor{tablehead}
    \tabhead{OE} & \tabhead{Factor (VI)} & \tabhead{Alternativas de diseño} \\
    \rowcolor{tablealt} \tabphase{OE1} & Representación intermedia & Paths/schemas vs behaviors/BO \\
    \rowcolor{white} \tabphase{OE2} & Schema matching & Correspondencia estructural + umbral \\
    \rowcolor{tablealt} \tabphase{OE3} & Similitud semántica & Embeddings / ontología / híbrido \\
    \rowcolor{white} \tabphase{OE4} & Pesos $\alpha$, $\beta$, $\gamma$ & $\alpha+\beta+\gamma=1$; tipologías \\
  \end{tabular}}
  \note{[9:15--9:55] Alex. VI = factores de diseño del investigador.}
\end{frame}

% --- 14 VD ---
\begin{frame}{Variables dependientes --- ratios por OE (Tabla 5 · Anexo D)}
  \beamertable{%
  \begin{tabular}{@{}L{0.07\linewidth}L{0.30\linewidth}L{0.56\linewidth}@{}}
    \rowcolor{tablehead}
    \tabhead{OE} & \tabhead{Ratio (VD)} & \tabhead{Rango / fórmula} \\
    \rowcolor{tablealt} \tabphase{OE1} & Completitud normalización & Extraídos / esperados \\
    \rowcolor{white} \tabphase{OE2} & StructuralScore (E) & $E \in [0, 1]$ \\
    \rowcolor{tablealt} \tabphase{OE3} & SemanticScore (S) & $S \in [0, 1]$ \\
    \rowcolor{white} \tabphase{OE4} & C; AlignmentScore & $C \in [0,1]$; $\alpha S+\beta E+\gamma C$ \\
  \end{tabular}}
  \note{[9:55--10:35] Alex. VD = AlignmentScore + clasificación.}
\end{frame}

% --- 15 Metodología ---
\begin{frame}{Metodología --- hitos / fases con verbos (§4.2 · §4.2.1)}
  \beamertable{%
  \resizebox{\linewidth}{!}{%
  \begin{tabular}{@{}L{1.05cm}L{1.85cm}L{4.55cm}L{4.05cm}@{}}
    \rowcolor{tablehead}
    \tabhead{Fase} & \tabhead{Verbo} & \tabhead{Transformación} & \tabhead{Producto} \\
    \rowcolor{tablealt} \tabphase{1} & Diseñar & Entradas, scores, umbrales & Modelo S/E/C documentado \\
    \rowcolor{white} \tabphase{2} & Normalizar & Contrato + SD (OE1) & Representaciones comparables \\
    \rowcolor{tablealt} \tabphase{3} & Emparejar & Estructuras (OE2) & StructuralScore E \\
    \rowcolor{white} \tabphase{4--5} & Calcular / Integrar & S (OE3) · C + score (OE4) & AlignmentScore + tipologías \\
  \end{tabular}}}
  \note{[10:35--11:15] Alex. Partir por metodología; en slide 16 ver modelo S/E/C (PI-14).}
\end{frame}

% --- 16 Modelo S/E/C ---
\begin{frame}{Artefacto --- modelo de alineación S / E / C}
  \figplan{slide17_img01}
  \figcaption{AlignmentScore $= \alpha\cdot S + \beta\cdot E + \gamma\cdot C$ $\rightarrow$ Alta / Media / Baja / Nula}
  \note{[11:15--12:00] Alex. Shvaiko (2005) sustenta E; scores por UA, agregado por contrato$\times$SD.}
\end{frame}

% --- 17 Matriz de consistencia (Anexo C) ---
\begin{frame}[t]{Matriz de consistencia --- Anexo C}
  \vspace{-0.85em}
  \tablesectiontitle{Problema tecnológico $\rightarrow$ objetivo $\rightarrow$ producto}
  \beamertablecompact{%
  \begin{tabularx}{\linewidth}{@{}L{0.05\linewidth}YYY@{}}
    \rowcolor{tablehead}
    \tabhead{F} & \tabhead{Problema tecnológico} & \tabhead{Objetivo} & \tabhead{Producto verificable} \\
    \rowcolor{tablealt} \tabphase{1} & Método insuficiente OpenAPI$\leftrightarrow$BIAN & Diseñar S/E/C & Modelo metodológico (Anexo C) \\
    \rowcolor{white} \tabphase{2} & Representaciones no comparables & OE1 Normalizar & Representaciones versionadas \\
    \rowcolor{tablealt} \tabphase{3} & Sin StructuralScore (E) & OE2 Emparejar & Matriz + $E \in [0,1]$ \\
    \rowcolor{white} \tabphase{4} & Sin SemanticScore (S) & OE3 Similitud & Matriz + $S \in [0,1]$ \\
    \rowcolor{tablealt} \tabphase{5} & Sin score integrado & OE4 Integrar & $C$; AlignmentScore; clasif. \\
  \end{tabularx}}
  \tablesectiontitle{VI · VD · Hipótesis}
  \beamertablecompact{%
  \begin{tabularx}{\linewidth}{@{}L{0.05\linewidth}YYY@{}}
    \rowcolor{tablehead}
    \tabhead{F} & \tabhead{VI} & \tabhead{VD} & \tabhead{Hipótesis} \\
    \rowcolor{tablealt} \tabphase{1} & Esquema S/E/C; $\alpha,\beta,\gamma$ & Definición scores & H0: modelo reproducible \\
    \rowcolor{white} \tabphase{2} & Esquema intermedio & Completitud OE1 & H1 (OE1) \\
    \rowcolor{tablealt} \tabphase{3} & Schema matching & StructuralScore E & H2 (OE2) \\
    \rowcolor{white} \tabphase{4} & Similitud semántica & SemanticScore S & H3 (OE3) \\
    \rowcolor{tablealt} \tabphase{5} & Pesos $\alpha,\beta,\gamma$ & C; AlignmentScore & H4 (OE4) \\
  \end{tabularx}}
  \note{[12:00--13:00] Alex. Matriz $\neq$ brecha Cap.\ 3. Señalar columnas; Anexo C completo en PDF del plan.}
\end{frame}

% --- 18 Brecha ---
\begin{frame}{Brecha de investigación V1--V5 (§3.6) --- no confundir con matriz}
  \figplan{Figura_06_Brecha}
  \figcaption{PI-11: Matriz Anexo C = problema tecnológico del diseño · Brecha V1--V5 = vacíos del estado del arte}
  \note{[13:20--14:00] Jack: V1--V2 · Alex: V3--V5.}
\end{frame}

% --- 19 Conclusiones ---
\begin{frame}{Conclusiones --- relaciones entre conceptos}
  \begin{planpanel}
  \begin{itemize}
    \item \keyterm{Objeto} (contrato) + \keyterm{referencia} (BIAN) + \keyterm{UA} $\rightarrow$ scores S/E/C integrados
    \item OE1--OE4 encadenan transformaciones con productos verificables
    \item Matriz Anexo C alinea problema tecnológico, objetivos, VI/VD e hipótesis
    \item Guía 1 entrega diseño; Tesis 2: artefacto software y validación empírica
  \end{itemize}
  \end{planpanel}
  \note{[14:00--14:30] Jack. Concluir con relaciones, no opinión (ppt-rules).}
\end{frame}

% --- 20 Recomendaciones ---
\begin{frame}{Recomendaciones}
  \begin{planpanel}
  \begin{itemize}
    \item Implementar artefacto software según procedimiento §4.2 (Tesis 2)
    \item Calibrar umbrales 0{,}80/0{,}60/0{,}40 con muestra bancaria representativa
    \item Validar hipótesis H1--H4 con contratos OpenAPI 3.x públicos (Anexo F)
  \end{itemize}
  \end{planpanel}
  \note{[14:30--14:45] Alex. Breve; continuidad post-Guía 1.}
\end{frame}

% --- 21 Referencias ---
\begin{frame}{Referencias citadas en esta exposición}
  \begin{planpanel}
  \begin{itemize}
    \item Casas, L., Pinto, S.\ \& Silva, J.\ (2021). OpenAPI specification survey. \textit{J. Syst. Softw.}
    \item Shvaiko, P.\ \& Euzenat, J.\ (2005). A survey of schema-based matching approaches. \textit{J. Data Semant.}
    \item Farzi, H.\ et al.\ (2021). BIAN service landscape. BIAN documentation.
  \end{itemize}
  \end{planpanel}
  \vspace{0.35em}
  {\large Otras referencias sectoriales (CMA, TSB): ver informe completo.}
  \note{[14:45--15:00] Jack. PDFs Casas y Shvaiko abiertos. ¿Preguntas?}
\end{frame}

\end{document}
